\section{Introducción: la era de la modelación generativa}

Estamos en una época de enorme transformación en la IA generativa (Generative AI). El objetivo central de la modelación generativa no es solo identificar patrones en los datos y tomar decisiones con base en ellos, sino, más importante aún, poder \textbf{generar instancias de datos completamente nuevas a partir de los patrones aprendidos}. Esta idea aparentemente simple ha provocado en los últimos años cambios revolucionarios en numerosos campos, entre ellos la generación de imágenes, la generación de audio y el procesamiento de lenguaje natural.

\begin{figure}[H]
\centering
\includegraphics[width=0.85\textwidth]{slides-images/slide-02.jpg}
\caption{¿Cuál de estos rostros es real? En realidad, los tres son rostros falsos generados por IA\protect\footnotemark}
\end{figure}
\footnotetext{Slides página 2. Fuente: thispersondoesnotexist.com}

La clase comienza con un ejemplo que invita a la reflexión: se muestran tres fotografías de rostros y se le pide a la audiencia que determine cuál es real. Al revelarse la respuesta, resulta que los tres fueron sintetizados por un modelo generativo. Esto ilustra de manera contundente la potencia de los modelos generativos modernos y da pie al tema central de esta sesión: comprender y dominar los fundamentos teóricos de la modelación generativa profunda.

\begin{importantbox}{Objetivo central de la modelación generativa}
Dadas muestras de entrenamiento provenientes de cierta distribución, aprender un modelo $P_{\text{model}}(x)$ capaz de representar esa distribución, de modo que se acerque lo más posible a la distribución real de los datos $P_{\text{data}}(x)$. Los modelos generativos pueden usarse en dos grandes tareas:
\begin{itemize}
    \item \textbf{Estimación de densidad} (Density Estimation): aprender la distribución de probabilidad subyacente de los datos
    \item \textbf{Generación de muestras} (Sample Generation): muestrear a partir de la distribución aprendida para generar datos nuevos
\end{itemize}
\end{importantbox}

\subsection{Aprendizaje supervisado vs. aprendizaje no supervisado}

En el contenido previo del curso nos hemos enfocado principalmente en problemas de \textbf{aprendizaje supervisado} (Supervised Learning): dados datos etiquetados $(x, y)$, el objetivo es aprender la función que mapea la entrada $x$ a la etiqueta $y$. La modelación generativa, en cambio, pertenece al ámbito del \textbf{aprendizaje no supervisado} (Unsupervised Learning): solo contamos con datos $x$, sin etiquetas, y el objetivo es descubrir la estructura subyacente oculta en los datos.

\begin{figure}[H]
\centering
\includegraphics[width=0.85\textwidth]{slides-images/slide-04.jpg}
\caption{Comparación entre aprendizaje supervisado y aprendizaje no supervisado\protect\footnotemark}
\end{figure}
\footnotetext{Slides página 4.}

\begin{table}[H]
\centering
\caption{Comparación entre aprendizaje supervisado y aprendizaje no supervisado}
\begin{tabular}{lll}
\toprule
\textbf{Dimensión} & \textbf{Aprendizaje supervisado} & \textbf{Aprendizaje no supervisado} \\
\midrule
Datos & $(x, y)$, $x$ son los datos, $y$ es la etiqueta & Solo $x$, sin etiquetas \\
Objetivo & Aprender la función de mapeo $x \to y$ & Descubrir la estructura oculta/subyacente de los datos \\
Tareas típicas & Clasificación, regresión, detección de objetos & Agrupamiento, extracción de características, reducción de dimensionalidad \\
\bottomrule
\end{tabular}
\end{table}

\subsection{Principio central de la modelación generativa}

\begin{figure}[H]
\centering
\includegraphics[width=0.85\textwidth]{slides-images/slide-05.jpg}
\caption{Modelación generativa: de la estimación de densidad a la generación de muestras\protect\footnotemark}
\end{figure}
\footnotetext{Slides página 5.}

El problema fundamental de la modelación generativa puede formularse así: ¿cómo aprender una distribución de probabilidad $P_{\text{model}}(x)$ que se acerque lo más posible a la distribución real de los datos $P_{\text{data}}(x)$?

$$P_{\text{model}}(x) \approx P_{\text{data}}(x)$$

Una vez que hemos logrado aprender un modelo de distribución de probabilidad de este tipo, podemos muestrear a partir de él para generar instancias de datos completamente nuevas, nunca antes vistas. Las muestras generadas reflejarán las características de la distribución de los datos de entrenamiento.

\subsection{¿Por qué nos interesan los modelos generativos?}

Los modelos generativos tienen tres grandes escenarios de aplicación en el mundo real:

\textbf{1. Reducción de sesgos (Debiasing)}: los modelos generativos pueden descubrir automáticamente la distribución subyacente de las características en un conjunto de datos. Por ejemplo, ante un conjunto de datos de rostros, el modelo puede detectar si características como el tono de piel, la pose o la iluminación están sobrerrepresentadas o subrepresentadas, lo que ayuda a crear conjuntos de datos más justos y representativos.

\begin{figure}[H]
\centering
\includegraphics[width=0.85\textwidth]{slides-images/slide-06.jpg}
\caption{Uso de modelos generativos para descubrir sesgos en conjuntos de datos\protect\footnotemark}
\end{figure}
\footnotetext{Slides página 6. Fuente: Amini/Soleymani+ AAAI/AIES 2019.}

\textbf{2. Detección de anomalías (Outlier Detection)}: en escenarios de seguridad crítica como la conducción autónoma, los modelos generativos pueden aprender la distribución de probabilidad de los datos para identificar eventos poco frecuentes situados en la cola de la distribución (como condiciones climáticas extremas o peatones), lo que permite detectar estas anomalías y ajustar el comportamiento del sistema.

\begin{figure}[H]
\centering
\includegraphics[width=0.85\textwidth]{slides-images/slide-07.jpg}
\caption{Modelos generativos aplicados a la detección de anomalías en la conducción autónoma\protect\footnotemark}
\end{figure}
\footnotetext{Slides página 7. Fuente: Amini+ IROS 2018.}

\textbf{3. Generación de muestras (Sample Generation)}: esta es la aplicación más conocida: usar modelos generativos para crear instancias de datos completamente nuevas. Esta es la base de aplicaciones de Generative AI como ChatGPT, la generación de imágenes y la generación de video.

\begin{figure}[H]
\centering
\includegraphics[width=0.85\textwidth]{slides-images/slide-08.jpg}
\caption{Diversas aplicaciones impulsadas por modelos generativos: lenguaje natural, imagen y video, biología\protect\footnotemark}
\end{figure}
\footnotetext{Slides página 8.}

\subsection{Resumen de la sección}

En esta sección se introdujeron los conceptos básicos de la modelación generativa y se explicó la transición del aprendizaje supervisado al aprendizaje no supervisado. El objetivo central de los modelos generativos es aprender la distribución de probabilidad de los datos, y pueden usarse en escenarios como la estimación de densidad, la detección de anomalías y la generación de muestras. Esta clase se centrará en dos arquitecturas fundamentales: el \textbf{Variational Autoencoder (VAE)} y la \textbf{Generative Adversarial Network (GAN)}, ambas pertenecientes a los \textbf{Latent Variable Model} (modelos de variables latentes).

